\usepackage[spanish]{babel} 
\usepackage[utf8]{inputenc}
%\usepackage[T1]{fontenc} %Codificacion de fuente
%\usepackage[latin1]{inputenc} %Codificacion de fuente
%\usepackage{cite} %% No compatible con Biblatex
\usepackage{comment}
\usepackage{graphicx}
\usepackage{listings}
\usepackage{amsmath,amssymb,amsfonts}
\usepackage{algorithmic}
\usepackage{textcomp}
\usepackage[table,xcdraw]{xcolor}
\usepackage{longtable}
\usepackage{threeparttable}
\usepackage{url}
\usepackage{booktabs}
\usepackage{float}
\usepackage{hyperref} % \usepackage[hidelinks]{hyperref}
\usepackage{enumitem}
\usepackage{multirow,multicol}
\usepackage{karnaugh-map}
\usepackage{csquotes}
\usepackage{marvosym}
\usepackage[normalem]{ulem}
\usepackage{array}
\usepackage{cleveref}
\usepackage{lineno}
\usepackage{ragged2e}
\usepackage{setspace}
\usepackage{pgffor}
\usepackage{etoolbox}
\usepackage{tikz}
%\usepackage[dvips,final]{epsfig}
\usepackage{lipsum}


\usetikzlibrary{matrix,calc}


\hyphenation{in-te-rrup-to-res}


\hypersetup{
    colorlinks=false,
    hyperindex=true,
    breaklinks=true,
    %linkcolor=blue,
    %filecolor=magenta,      
    %urlcolor=cyan,
    %citecolor=blue,
    pdftitle={Trabajo UNAD},
    pdfpagemode=FullScreen,
}


\selectlanguage{spanish}

% Colores estilo "Dark Mode" suave o editor moderno
\definecolor{codegreen}{rgb}{0,0.6,0}
\definecolor{codegray}{rgb}{0.5,0.5,0.5}
\definecolor{codepurple}{rgb}{0.58,0,0.82}
\definecolor{arduino_teal}{rgb}{0.0, 0.59, 0.62}
\definecolor{backcolour}{rgb}{0.95,0.95,0.92}

% ESTILO BASE (Con soporte de tildes incluido)
\lstdefinestyle{base_style}{
    backgroundcolor=\color{backcolour},   
    numberstyle=\tiny\color{codegray},
    basicstyle=\footnotesize\ttfamily,
    breaklines=true,                 
    numbers=left,                    
    numbersep=7pt,
    xleftmargin=10pt,                  
    showstringspaces=false,
    frame=single,
    rulecolor=\color{black!10},
    literate={á}{{\'a}}1 {é}{{\'e}}1 {í}{{\'i}}1 {ó}{{\'o}}1 {ú}{{\'u}}1
             {Á}{{\'A}}1 {É}{{\'E}}1 {Í}{{\'I}}1 {Ó}{{\'O}}1 {Ú}{{\'U}}1
             {ñ}{{\~n}}1 {Ñ}{{\~N}}1
             {α}{{\ensuremath{\alpha}}}1 {β}{{\ensuremath{\beta}}}1 
             {γ}{{\ensuremath{\gamma}}}1 {δ}{{\ensuremath{\delta}}}1
             {λ}{{\ensuremath{\lambda}}}1 {π}{{\ensuremath{\pi}}}1
             {Ω}{{\ensuremath{\Omega}}}1 {τ}{{\ensuremath{\tau}}}1
	% Soporte español
}

% ESTILOS ESPECÍFICOS
\lstdefinestyle{python_style}{
    style=base_style,
    language=Python,
    commentstyle=\color{codegreen},
    keywordstyle=\color{magenta},
    stringstyle=\color{codepurple}
}

\lstdefinestyle{arduino_style}{
    style=base_style,
    language=C++,
    commentstyle=\color{gray},
    keywordstyle=\color{arduino_teal}\bfseries,
    stringstyle=\color{blue}
}

\lstdefinestyle{bash_style}{
    style=base_style,
    language=bash,
    basicstyle=\ttfamily\footnotesize\color{white},
    backgroundcolor=\color{black!85},
    keywordstyle=\color{cyan},
    numbers=none
}


\providecommand{\inputpath}{./src/} % Carpeta sugerida para tus códigos
\DeclareGraphicsExtensions{.pdf,.jpeg,.png,.eps}

\graphicspath{ {./graphics/} {./img/} {./diagramas/} }

%%%%%%%%%%%%%%%%%%%%%%%%%%%%%%%%%%%%%%%%%%%%%%%%%%%%%%%%%%%
% LÓGICA DINÁMICA APA vs IEEE
%\makeatletter
\ifx\estilo\apa
	%\renewcommand\contentsname{Índice}
	\addto\captionsspanish{\renewcommand{\listtablename}{Índice de tablas}}
	\renewcommand{\lstlistlistingname}{Índice de códigos}
	
	\newcommand{\insertarFigura}[4]{
        \begin{figure}[ht]
            \centering
            \caption{#3}
            \label{#4}
            \includegraphics[width=#2\linewidth]{#1}
        \end{figure}
    }
    \newcommand{\insertarFiguraConNota}[5]{
    		\begin{figure}[ht]
        		\centering
	        \caption{#3} 
	        \label{#4}
    		    \includegraphics[width=#2\linewidth]{#1}
	        \ifx\estilo\apa \par\vspace{2pt}\small \textit{Nota.} #5 \else \caption{#3} \label{#4} \fi
	    \end{figure}
	}
	
    %\lstset{captionpos=t} % APA 7: Título arriba
    %\AtBeginEnvironment{lstlisting}{\captionsetup{labelfont=bf, textfont=it}}

	\usepackage{caption}
	\DeclareCaptionLabelSeparator{tightnewline}{\\[-1.5ex]}
	
	%\captionsetup[lstlisting]{
	%    font=small,       % Tamaño típico para IEEE e APA
    %		labelfont=bf,     % "Listing 1" en negrita
	%    singlelinecheck=false,
    	%	margin=0pt
	%}
	\captionsetup[lstlisting]{position=top, labelfont=bf, textfont=it, singlelinecheck=false, margin=0pt, labelsep=tightnewline}
	%\usepackage{caption}
    %% 1. Definimos un formato de caption con interlineado APA (1.5 o 2.0)
    %\DeclareCaptionFormat{apa-listings}{
    %    \textbf{#1#2} \\[-0.3cm] %[0.3cm]
    %    \textit{#3}
    %}
    %\captionsetup[lstlisting]{
    %    format=apa-listings,
    %    font={stretch=1.0},
    %    skip=13pt,
    %    labelsep=none,
    %    singlelinecheck=false,
    %    justification=raggedright
    %}
    %\lstset{
    %    style=base_style,
    %    captionpos=t, 
    %    basicstyle=\ttfamily\small,
    %    lineskip=0pt,
    %    linewidth=1.0\linewidth,
    %    xleftmargin=0.04\linewidth,
    %}
\else
    \renewcommand{\IEEEkeywordsname}{Palabras Clave}
	\renewcommand{\bibfont}{\footnotesize}
	\setlength{\biblabelsep}{0.5em} % Espacio tras la etiqueta
	\setlength{\bibhang}{0pt}% Sangría para líneas siguientes
	\addto\captionsspanish{\renewcommand{\figurename}{Fig.}}
	
	\newcommand{\insertarFigura}[4]{
        \begin{figure}[ht]
            \centering
            \includegraphics[width=#2\linewidth]{#1}
            \caption{#3}
            \label{#4}
        \end{figure}
    }
    %\captionsetup[lstlisting]{position=bottom, labelfont=sc} % Small caps para IEEE

	% Si necesitas forzar el estilo de la fuente solo para la etiqueta del listing:
	\usepackage{etoolbox}
	\AtBeginEnvironment{lstlisting}{\captionsetup{labelfont=sc, textfont=it}}
    %\lstset{
    %    style=base_style,
    %    captionpos=b, 
    %    basicstyle=\ttfamily\footnotesize,
    %    lineskip=0pt,
    %    linewidth=1.0\linewidth,
    %    xleftmargin=0.05\linewidth,
    %}

\fi
%\makeatother

\ifx\estilo\apa
\else
    \lstset{captionpos=b} % IEEE: Título abajo
    \AtBeginEnvironment{lstlisting}{\captionsetup{labelfont=sc, textfont=rm}}
\fi

\addto\captionsspanish{\renewcommand{\tablename}{Tabla}}
\renewcommand{\lstlistingname}{Código}


% --- ENTORNO PARA OBJETIVOS DINÁMICOS ---
\newenvironment{misObjetivosEspecificos}{
    \ifx\estilo\apa
        % Si es APA, no hace nada especial (párrafos normales)
        \par\medskip
    \else
        % Si es IEEE, inicia una lista de puntos
        \begin{itemize}[leftmargin=*]
    \fi
}{
    \ifx\estilo\apa
        \par\medskip
    \else
        \end{itemize}
    \fi
}

% Comando para los ítems que se adapta
\newcommand{\objItem}{
    \ifx\estilo\apa 
        \par \space 
    \else 
        \item 
    \fi
}

\newcommand{\autoresFinales}{}
\newcommand{\thanksFinales}{Los autores pertenecen a la \miUniversidad, \miEscuela. Correos: }

\foreach \i in {1,...,\totalIntegrantes}{
    % 1. Extraer datos
    \expandafter\let\expandafter\tempNombre\csname nombre\i\endcsname
    \expandafter\let\expandafter\tempEmail\csname email\i\endcsname
    
    % 2. Acumular nombres para el encabezado (separados por comas)
    \xdef\tempData{\tempNombre\ifnum\i<\totalIntegrantes, \ \fi}
    \expandafter\gappto\expandafter\autoresFinales\expandafter{\tempData}
    
    % 3. Acumular en el Thanks Compuesto (Nombre <email>)
    \xdef\tempThanks{\tempNombre\ (\noexpand\texttt{\tempEmail})\ifnum\i<\totalIntegrantes; \else. \fi}
    \expandafter\gappto\expandafter\thanksFinales\expandafter{\tempThanks}
}
